\chapter{Problem Analysis and Requirements}
\label{cha:requirements}

The Norlys use case requires more than a plausible explanation of a failed
request. An operator needs to inspect the observations behind the diagnosis
and understand a proposed change before approving it. This chapter analyses
that need and derives requirements from the company specifications and
Chapter~\ref{cha:state-of-the-art}. The requirements connect the four research
questions to observable acceptance criteria. Chapter~5 will define how those
criteria are examined through experiments, while Chapters~6 and~7 will explain
the design and its implementation.

\section{Norlys Use Case and Project Boundary}
\label{sec:requirements:context}

The intended user is a Cloud Platform operations engineer investigating a
reported endpoint failure. In the workflow specified by Norlys, the operator
provides the endpoint and a time window. The system investigates observations
from Cloudflare, Azure Firewall, Azure infrastructure and the application,
then presents the request path, a proposed cause, supporting evidence and an
actionable repair. An approved repair can subsequently be applied through
the Cybertron portal \parencite[User Story, AC-4--AC-6]{norlys_operations_agents}.
The operational problem is the work needed to connect observations across
these layers. An application error alone may describe a symptom, while a
request blocked earlier in the path may leave no application record.

This investigation is a bounded part of Norlys' wider observability
initiative. The parent specification covers collection, normalisation,
protection of sensitive data, retention and access across a larger set of
operational sources \parencite{norlys_observability_spec}. The agent PoC should
consume that platform's prepared telemetry where it is available. Direct
source queries are permitted by the PoC specification where the platform
does not yet provide it, but this dependency gap must be documented
\parencite[Relationship to the Parent Story]{norlys_operations_agents}.
Building the complete observability platform is therefore outside this
thesis. Protecting the data used by the PoC remains a requirement even when
a connector queries a source directly.

The experimental boundary is a designated non-production application request
path. Authorised reading of operational telemetry is distinct from permission
to modify the resources that produced it. Production firewall records can
provide context, but their availability does not authorise fault injection
or remediation on a production firewall. Both activities must remain within
the permitted test environment. The source specification names a landing zone
\texttt{ai-operations}. The agreed project environment is instead
\texttt{cpCentralOperation-D}. This is an explicit adaptation of the resource
location, retaining the requirement for separation from customer workloads
and for restricted write permissions.

Within this boundary, security-related failures concern controls that block
legitimate traffic, such as an incorrect firewall or web application firewall
rule. Attack detection, production service operation, additional telemetry
sources and model training are excluded. As justified in
Section~\ref{sec:state-of-the-art:ml-llm}, developing and benchmarking a
separate conventional ML diagnosis model is also outside scope. Statistical
summaries of observations can still support the investigation.

Portal integration and human-approved remediation remain requirements from
the Norlys specification. A local console is an intermediate development
interface and does not satisfy the final portal requirement. Similarly, a
successful model connection establishes access, not diagnostic capability.
The source document's sketch of an orchestrator and one specialist per
telemetry source is a proposed architecture
\parencite[Architecture Sketch]{norlys_operations_agents}. The multi-agent objective is
retained, but agent roles, runtime placement and coordination mechanisms
require justification in System Design.

\thesisTODO{Confirm the designated Cloudflare-to-application test route and
the non-production resources on which faults and repairs are permitted.
Separate verified links from assumed links. Observed access through Global
Secure Access does not establish the complete Cloudflare route. Use this
inventory to identify any infeasible parts of the Norlys scope.}

\section{Analysis of the Main Challenges}
\label{sec:requirements:analysis}

\noindent\textbf{Connecting observations without overstating certainty.}
A shared request or operation identifier can link records directly. Other
sources may provide only times, addresses, ports and rule decisions. Several
requests can then fit the same observations. As discussed in
Section~\ref{sec:state-of-the-art:correlation}, reconstructing a path depends
on the available identifiers and on knowledge of the components and their
connections. The system must preserve the basis of each link. A relationship
inferred from timing and destination cannot be presented as an exact request
match. Missing records also need an explicit explanation where possible,
because a blocked request, an unavailable source and an empty query result
have different meanings.

\noindent\textbf{Distinguishing a relevant security event from background traffic.}
A denied connection near an incident is a candidate observation, not by
itself the cause. The investigation needs the action, affected destination,
rule and timing, together with evidence that the traffic belongs to the
selected path. A rule blocking a legitimate request should be distinguished
from an unrelated denial or an application returning an error after traffic
was allowed. Section~\ref{sec:state-of-the-art:security} motivates retaining
this context when selecting fields and calculating derived features. For
example, a count of denied connections is interpretable only with its time
window and filters. Whether such a count adds diagnostic value beyond the
individual rule event remains an evaluation question for RQ2.

\noindent\textbf{Producing a diagnosis that can be checked.}
Identifying the application as the affected layer does not necessarily
identify its failing dependency or configuration. Likewise, a correct cause
can be supported by irrelevant or invented evidence. The distinctions in
Section~\ref{sec:state-of-the-art:evaluation} require separate outputs for
the component, proposed cause and supporting observations. The operator must
also be able to recognise a healthy endpoint and an inconclusive
investigation. Absence of evidence is not sufficient to declare health, and
agreement between agents is not evidence that their conclusion is correct.

The organisation of the investigation introduces a related uncertainty.
Section~\ref{sec:state-of-the-art:agents} does not establish a general
advantage for multiple agents under comparable conditions. If a collaborating
system receives additional tools, more evidence or a larger total budget,
any improvement cannot be attributed to collaboration alone. Requirements
must therefore support comparable inputs, outputs and resource accounting
without fixing the final agent architecture.

\noindent\textbf{Turning a diagnosis into a controlled intervention.}
A useful repair suggestion is not yet an authorised action, and a successful
management command is not proof of recovery. Section~\ref{sec:state-of-the-art:remediation}
motivates separate checks for the proposed change, its execution, the endpoint
outcome and restoration of prior state. These checks must be enforced by the
system rather than depend on the model following an instruction. The
operator needs enough information to approve a specific change, while the
execution identity must be unable to modify resources outside the test scope.

\section{Functional Requirements}
\label{sec:requirements:functional}

The following requirements translate the analysed problems into observable
system behaviour. Their identifiers are used in the acceptance table and
can subsequently link design components and evaluation results to the
original need. The Norlys specification groups its expectations into seven
acceptance criteria, numbered AC-1--AC-7. References to those criteria identify
the company basis for the requirements below.

\noindent\textbf{F1: Scoped investigation and telemetry access.}
The system shall accept an endpoint and time window, assign an investigation
identifier and retrieve relevant observations from all four specified source
groups. These comprise Cloudflare requests, WAF events and origin errors,
firewall network and application rule events, infrastructure health, metrics,
activity and diagnostic configuration, and application requests, logs and
traces \parencite[AC-3]{norlys_operations_agents}.
Each connector shall be independently testable. Queries shall be
bounded by time and incident scope, and their results shall distinguish no
matching observations from query failure, access denial and unavailable
telemetry. This implements Norlys AC-3 and AC-4 while making retrieval
limitations visible to the investigator.

\noindent\textbf{F2: Traceable correlation.}
The system shall preserve source and resource identifiers, event times and
available request or operation identifiers. Each proposed relationship along
the path shall identify its supporting records and whether it is confirmed
by a suitable identifier or inferred from other observations. Ambiguous and
unobserved segments shall remain visible. A partial path is an acceptable
representation of incomplete evidence, but must not be reported as complete.

\noindent\textbf{F3: Interpretable security evidence.}
The evidence available for security-related diagnosis shall include, where
recorded, time, source and destination, request target, response status,
policy action and rule identifier. Unavailable fields shall be reported
explicitly. Derived features shall state their definition, filters, time
window and supporting observations. Initial candidates include counts of
allowed and denied traffic for the selected destination. Their inclusion
must be justified by availability and diagnostic relevance, rather than by
assuming that features useful for attack detection transfer to configuration
faults.

\noindent\textbf{F4: Structured diagnosis with explicit uncertainty.}
An investigation shall return a structured result identifying the suspected
faulty layer or component, proposed cause, evidence references, uncertainty
and unresolved alternatives. It shall also support outcomes for no fault
found in the investigated scope and insufficient evidence. A confidence
indication shall explain its basis. A model-generated score shall not be
interpreted as a probability of correctness without validation. The output
format must allow component identification, cause identification and evidence
support to be assessed separately.

\noindent\textbf{F5: Reviewable remediation plan.}
A proposed repair shall name the target resource, exact change, supporting
diagnosis, expected effect, validation checks and rollback procedure. Before
approval, the system shall perform a dry run, meaning a check of the proposed
change without applying it, and display its result and limitations. An action
without a supported validation or rollback procedure shall remain a manual
recommendation rather than become executable. Norlys' proposed action
catalogue includes rule changes, compute restart or scaling, health-probe or
network-security-rule correction, and certificate replacement
\parencite[AC-5]{norlys_operations_agents}. This catalogue does not itself
grant permission to execute those actions.

\noindent\textbf{F6: Enforced approval and action limits.}
Execution shall require explicit approval by an authorised operator for the
specific plan. The system shall check the action allowlist, parameters,
target scope and required preconditions before making a change. Changing an
approved plan shall require new approval. Rejected, unapproved and
out-of-scope actions shall not execute. Deletion, scaling to zero and DNS
record removal are excluded from the remediation allowlist, as required by
Norlys AC-5. These controls apply independently of the model's confidence.

\noindent\textbf{F7: Verified recovery and rollback.}
The system shall record execution separately from recovery. After a change,
it shall check the endpoint and relevant configuration or telemetry against
the expected effect. It shall retain the prior state needed for the defined
rollback and expose a one-step rollback action to an authorised operator.
Unsuccessful remediation shall be reported and either rolled back within
the approved procedure or escalated. Rollback shall have its own verification
result. Restoration of prior configuration does not automatically imply
that every external effect has been reversed.

\noindent\textbf{F8: Operator interaction and investigation history.}
Cybertron shall provide investigation input, the reconstructed path, diagnosis,
evidence and repair plan. Where a plan is executable, the portal shall show
its dry-run result and provide approval and rejection controls. Execution,
recovery and rollback outcomes shall remain accessible with past
investigations. Viewing and approval shall use distinct Entra ID permissions,
following the existing portal navigation and Platform Management conventions
\parencite[AC-6]{norlys_operations_agents}.

\noindent\textbf{F9: Repeatable, independently checkable cases.}
The project shall provide a repeatable way to introduce and reset known
faults in the test environment, together with healthy cases. The reference
answer shall record the injected condition and its independently verified
effect. Scenario names and other test metadata that disclose the cause shall
be kept separate from agent-visible telemetry, while retaining realistic
operational symptoms. The evaluation shall support repeated runs and cases
with missing or ambiguous evidence. Norlys AC-7 also requires unannounced
faults introduced by someone other than the agent developer and at least
one exercised rollback. Chapter~5 will define the cases and assessment
procedure.

The source fault catalogue spans Cloudflare origin, WAF and DNS failures,
firewall rules and destinations, infrastructure health and configuration,
and application exceptions, dependencies and configuration
\parencite[AC-7]{norlys_operations_agents}. Its full coverage is a company
expectation against which the selected evaluation scope must be checked.
For example, a production firewall on the observed path cannot be made a
fault-injection target merely to complete that catalogue.

\thesisTODO{Agree the fault catalogue and executable remediation allowlist
against the permitted test resources. Record each reduction or deferral
from Norlys AC-5 and AC-7 with its reason and effect on the RQs. Retain portal
integration and remediation as requirements unless a scope change is
explicitly agreed.}

\section{Non-functional Requirements and Constraints}
\label{sec:requirements:constraints}

\noindent\textbf{N1: Isolation and least privilege.}
The PoC shall use the agreed landing zone and keep investigation access
read-only. A separate remediation identity shall hold only the permissions
needed for approved changes to the non-production test resources, with no
production write permissions. Azure identity-based access shall be used
where supported. The specification also explicitly permits a read-only
Cloudflare API token, so its general preference for secret-free access cannot
be assumed to cover every connector. Necessary credentials must be scoped
and protected, and must not enter prompts, reports or source control
\parencite[AC-1, AC-3]{norlys_operations_agents}.

\noindent\textbf{N2: Controlled disclosure and retention.}
Queries, model inputs and stored investigation records shall contain only
data needed for the authorised investigation. Sensitive fields shall be
masked or omitted according to the agreed Norlys policy, and access and
retention shall cover retrieved evidence as well as conclusions. A direct
connector must apply the required protection when the parent platform has
not already done so. Resource ownership, resource region and the permitted
location of model processing are separate constraints.

\thesisTODO{Confirm with Norlys which operational fields may be sent to the
model, the required masking, permitted processing locations and retention
and access rules for investigation records. Sweden Central is the preferred
resource location, but Sweden-only model processing has not been established
as a project requirement.}

\noindent\textbf{N3: Reproducibility and auditability.}
Agent code, prompts, tool definitions and configuration shall be versioned.
Infrastructure shall be represented as code and deployed through the agreed
pipeline, as required by Norlys AC-1 and AC-2. Each investigation shall record
the relevant versions, inputs, queries, permitted evidence, tool outcomes
and final output. Remediation records shall additionally identify the
approver, approved change, before and after state, execution and verification
results. These records must permit inspection of how a conclusion was
reached. Repeating a recorded configuration is not a guarantee of identical
model output.

\noindent\textbf{N4: Bounded operation and comparable resource accounting.}
An investigation shall have configurable stopping limits and record elapsed
time, model calls, input and output tokens, and tool calls. Failures, retries,
timeouts and invalid outputs shall remain in the record. Estimated cost per
investigation shall state the usage and pricing assumptions behind it.
Limits shall apply to the entire investigation, including all collaborating
agents. The system
shall support a common evidence interface and output format so that candidate
organisations can be evaluated using the same model version, tools, cases
and available evidence under the same total resource limits. The proposed
single-investigator and collaborating comparison still requires a protocol
in Chapter~5. No minimum accuracy or
assumed collaboration advantage is specified here.

\noindent\textbf{N5: Model access and capability.}
The model deployment shall be controlled within the agreed Norlys landing
zone rather than depend on a model in another Norlys subscription. The
selected model and its integration shall support tool-directed investigation
and diagnostic output that can be checked against a defined structure.
Invalid tool requests or outputs must be detected and recorded. The model
version, deployment settings and applicable usage limits shall be documented
for the evaluation. These are selection criteria, not a conclusion that the
initial connectivity-test model is the best diagnostic model. Chapter~6 will
justify the selection against the criteria and any pilot evidence.

\section{Acceptance Criteria and Traceability}
\label{sec:requirements:acceptance}

Table~\ref{tab:requirements:traceability} connects each requirement to its
source, research question and observable acceptance evidence. References to
Norlys AC-1--AC-7 denote the numbered acceptance criteria in the
\emph{Operations with AI Agents} specification. References to Chapter~3
identify the research finding that motivates the requirement.

Acceptance of a capability and effectiveness of the resulting system are
different questions. For example, the presence of evidence references can
be checked as a software property, but whether the cited observations support
the diagnosis requires assessment against the case. Likewise, preventing an
unapproved action does not establish that an approved repair is effective.
The table identifies observable checks. Diagnostic and remediation
effectiveness remain questions for the evaluation.

Chapter~5 will specify scenario selection, repetitions, scoring rules,
comparison conditions and budgets. It will assess component and cause
correctness, evidence support, healthy-case behaviour, consistency and
resource use separately. Remediation assessment will distinguish plan
quality, approval enforcement, execution, recovery and rollback. A result
showing no advantage from collaboration can still answer the research
questions. The implemented requirements, deviations and observed outcomes
must remain traceable through System Design, Implementation and Evaluation
and Results.

\begin{landscape}
\begingroup
\fontsize{9.5}{11}\selectfont
\setlength{\tabcolsep}{4pt}
\renewcommand{\arraystretch}{1.0}
\begin{longtable}{@{}>{\raggedright\arraybackslash}p{0.045\linewidth}>{\raggedright\arraybackslash}p{0.14\linewidth}>{\raggedright\arraybackslash}p{0.19\linewidth}>{\raggedright\arraybackslash}p{0.065\linewidth}>{\raggedright\arraybackslash}p{0.28\linewidth}>{\raggedright\arraybackslash}p{\dimexpr0.28\linewidth-10\tabcolsep\relax}@{}}
\caption{Requirements, rationale and acceptance evidence.}
\label{tab:requirements:traceability}\\
\toprule
ID & Requirement & Source/rationale & Relevant RQ & Acceptance criterion & Verification method \\
\midrule
\endfirsthead
\multicolumn{6}{@{}l}{\tablename~\thetable\ continued}\\
\toprule
ID & Requirement & Source/rationale & Relevant RQ & Acceptance criterion & Verification method \\
\midrule
\endhead
\midrule
\multicolumn{6}{r@{}}{Continued on next page}\\
\endfoot
\bottomrule
\endlastfoot
F1 & Scoped input and retrieval & Norlys AC-3, AC-4. Bounded source access. & RQ1 & All four connectors use endpoint and time scope, distinguishing records, no matches and failures. & Test each connector and combined retrieval. Inspect query scope and error outcomes. \\
\addlinespace
F2 & Traceable correlation & Section~\ref{sec:state-of-the-art:correlation}. Preserve link uncertainty. & RQ1, RQ3 & Links retain evidence and matching basis. Missing and ambiguous segments remain visible. & Check known links, missing records and competing requests against case records. \\
\addlinespace
F3 & Security fields and features & Section~\ref{sec:state-of-the-art:security}, Norlys AC-3. Rule decisions need context. & RQ2 & Field availability is documented. Aggregates have definitions, scope and traceable inputs. & Inspect rule events. Recalculate aggregates and assess their diagnostic contribution. \\
\addlinespace
F4 & Checkable diagnosis & Norlys AC-4, Section~\ref{sec:state-of-the-art:evaluation}. Assess correctness and grounding. & RQ3 & Output separates component, cause, evidence and uncertainty, including healthy and inconclusive outcomes. & Validate output and references. Score component, cause and evidence against known cases. \\
\addlinespace
F5 & Repair plan and dry run & Norlys AC-5, Section~\ref{sec:state-of-the-art:remediation}. Make proposed changes reviewable. & RQ4 & Plan states target, change, expected effect, dry-run result and rollback before approval. & Inspect plans, including actions lacking supported validation or rollback. \\
\addlinespace
F6 & Approval and allowlist & Norlys AC-5, Section~\ref{sec:state-of-the-art:remediation}. Enforce permission separately. & RQ4 & Unapproved, rejected, changed and out-of-scope plans cannot execute. Forbidden actions are rejected. & Test permitted and forbidden transitions. Inspect audit records and resource state. \\
\addlinespace
F7 & Recovery and rollback & Norlys AC-5, AC-7, Section~\ref{sec:state-of-the-art:remediation}. Verify state after action. & RQ4 & Execution, recovery and rollback have separate results. Permitted changes support restoration of prior state. & Check endpoint and state after successful and failed repairs and at least one rollback. \\
\addlinespace
F8 & Portal and history & Norlys AC-6. Operator visibility and role separation. & RQ3, RQ4 & Cybertron supports input, review, approval/rejection, outcomes, rollback and history with distinct permissions. & Walk through the portal using viewing and approval roles, including rejection. \\
\addlinespace
F9 & Repeatable controlled cases & Norlys AC-7, Section~\ref{sec:state-of-the-art:evaluation}. Separate symptoms and answer labels. & RQ1--4 & Cases can be reset. Reference answers stay hidden. Healthy and independently injected unannounced cases are included. & Inspect fault records and inputs. Repeat setup and reset. Verify injected effects independently. \\
\addlinespace
N1 & Isolation and least privilege & Norlys AC-1, AC-3, AC-5. Separate read and write authority. & RQ1, RQ4 & Investigation is read-only. Separate execution permissions exclude production writes. Credentials are protected. & Review permissions and credentials. Check write rejection on an out-of-scope test resource. \\
\addlinespace
N2 & Data protection and retention & Parent specification. Protection applies to direct queries too. & RQ1--4 & Fields, masking, processing locations, retention and access conform to the agreed policy. & Inspect query results, model inputs and stored records against the confirmed policy. \\
\addlinespace
N3 & Versioning and auditability & Norlys AC-1, AC-2, AC-5, Section~\ref{sec:state-of-the-art:evaluation}. Inspectable run records. & RQ1--4 & Records identify versions and evidence. Actions retain approver, change, before/after state and outcomes. & Trace conclusions to evidence and versions, and repairs through approval and verification. \\
\addlinespace
N4 & Limits and accounting & Norlys AC-2, Sections~\ref{sec:state-of-the-art:agents} and~\ref{sec:state-of-the-art:evaluation}. Comparable resources. & RQ3 & Limits cover all agents, retries and failures. Usage and cost assumptions are recorded. & Test stopping and failures. Reconcile usage and compare model, evidence, tools and budgets. \\
\addlinespace
N5 & Model access and capability & Agreed landing zone, Norlys AC-2, F1 and F4. Model selection constraints. & RQ1, RQ3 & Deployment is in the agreed landing zone. Tool requests and diagnostic outputs are validated, with versions recorded. & Inspect ownership and settings. Test valid and invalid tool and output formats. \\
\end{longtable}
\endgroup
\end{landscape}
